% Options for packages loaded elsewhere
\PassOptionsToPackage{unicode}{hyperref}
\PassOptionsToPackage{hyphens}{url}
%
\documentclass[
]{article}
\usepackage{amsmath,amssymb}
\usepackage{iftex}
\ifPDFTeX
  \usepackage[T1]{fontenc}
  \usepackage[utf8]{inputenc}
  \usepackage{textcomp} % provide euro and other symbols
\else % if luatex or xetex
  \usepackage{unicode-math} % this also loads fontspec
  \defaultfontfeatures{Scale=MatchLowercase}
  \defaultfontfeatures[\rmfamily]{Ligatures=TeX,Scale=1}
\fi
\usepackage{lmodern}
\ifPDFTeX\else
  % xetex/luatex font selection
\fi
% Use upquote if available, for straight quotes in verbatim environments
\IfFileExists{upquote.sty}{\usepackage{upquote}}{}
\IfFileExists{microtype.sty}{% use microtype if available
  \usepackage[]{microtype}
  \UseMicrotypeSet[protrusion]{basicmath} % disable protrusion for tt fonts
}{}
\makeatletter
\@ifundefined{KOMAClassName}{% if non-KOMA class
  \IfFileExists{parskip.sty}{%
    \usepackage{parskip}
  }{% else
    \setlength{\parindent}{0pt}
    \setlength{\parskip}{6pt plus 2pt minus 1pt}}
}{% if KOMA class
  \KOMAoptions{parskip=half}}
\makeatother
\usepackage{xcolor}
\usepackage{longtable,booktabs,array}
\usepackage{calc} % for calculating minipage widths
% Correct order of tables after \paragraph or \subparagraph
\usepackage{etoolbox}
\makeatletter
\patchcmd\longtable{\par}{\if@noskipsec\mbox{}\fi\par}{}{}
\makeatother
% Allow footnotes in longtable head/foot
\IfFileExists{footnotehyper.sty}{\usepackage{footnotehyper}}{\usepackage{footnote}}
\makesavenoteenv{longtable}
\setlength{\emergencystretch}{3em} % prevent overfull lines
\providecommand{\tightlist}{%
  \setlength{\itemsep}{0pt}\setlength{\parskip}{0pt}}
\setcounter{secnumdepth}{-\maxdimen} % remove section numbering
\ifLuaTeX
  \usepackage{selnolig}  % disable illegal ligatures
\fi
\IfFileExists{bookmark.sty}{\usepackage{bookmark}}{\usepackage{hyperref}}
\IfFileExists{xurl.sty}{\usepackage{xurl}}{} % add URL line breaks if available
\urlstyle{same}
\hypersetup{
  hidelinks,
  pdfcreator={LaTeX via pandoc}}

\author{}
\date{}

\begin{document}

\hypertarget{ripasso-prerequisiti}{%
\section{Ripasso prerequisiti}\label{ripasso-prerequisiti}}

\hypertarget{le-funzioni}{%
\subsection{1. LE FUNZIONI}\label{le-funzioni}}

Immagina una~\textbf{macchinetta mangiasoldi}:

\begin{itemize}
\tightlist
\item
  Tu metti una moneta (input)
\item
  La macchina ti dà una lattina (output)
\item
  Ogni moneta → una specifica lattina
\end{itemize}

Una~\textbf{funzione matematica}~è esattamente così:

\begin{itemize}
\tightlist
\item
  È una~\textbf{``macchina''}~che trasforma numeri
\item
  Tu gli dai un numero in~\textbf{ingresso}~(chiamato~\texttt{x})
\item
  Lei ti restituisce un numero in~\textbf{uscita}~(chiamato~\texttt{f(x)}~o~\texttt{y})
\end{itemize}

Quando scrivo:

\begin{verbatim}
f(x) = 2x + 3
\end{verbatim}

Significa:

\begin{itemize}
\tightlist
\item
  \texttt{f}~= nome della macchina/funzione
\item
  \texttt{(x)}~= il numero che tu metti nella macchina
\item
  \texttt{2x\ +\ 3}~= cosa fa la macchina con il numero
\end{itemize}

\textbf{Esempio pratico:}

Se~\texttt{x\ =\ 4}:

\begin{enumerate}
\def\labelenumi{\arabic{enumi}.}
\tightlist
\item
  Prendi~\texttt{x\ =\ 4}
\item
  La macchina fa:~\texttt{2\ ×\ 4\ +\ 3}
\item
  Risultato:~\texttt{8\ +\ 3\ =\ 11}
\item
  Quindi~\texttt{f(4)\ =\ 11}
\end{enumerate}

Se ho una funzione:~\texttt{g(t)\ =\ t²\ -\ 1}

Cosa significa~\texttt{g(3)}? → 3² - 1 = 9 - 1 = 8

\hypertarget{le-onde}{%
\subsection{2. LE ONDE}\label{le-onde}}

\hypertarget{immagina-di-essere-in-spiaggia}{%
\subsubsection{\texorpdfstring{\textbf{Immagina di essere in spiaggia:}}{Immagina di essere in spiaggia:}}\label{immagina-di-essere-in-spiaggia}}

\begin{enumerate}
\def\labelenumi{\arabic{enumi}.}
\tightlist
\item
  \textbf{L'acqua sale}~(arriva verso la spiaggia)
\item
  \textbf{Raggiunge il punto più alto}~(cresta dell'onda)
\item
  \textbf{L'acqua scende}~(torna indietro)
\item
  \textbf{Raggiunge il punto più basso}~(ventre dell'onda)
\item
  \textbf{Poi ricomincia da capo}
\end{enumerate}

Questo movimento~\textbf{SU e GIÙ}~che~\textbf{SI RIPETE}~è un'\textbf{onda}.

\hypertarget{altri-esempi-concreti-di-onde-nella-vita}{%
\subsubsection{\texorpdfstring{\textbf{ALTRI ESEMPI CONCRETI DI ONDE NELLA VITA:}}{ALTRI ESEMPI CONCRETI DI ONDE NELLA VITA:}}\label{altri-esempi-concreti-di-onde-nella-vita}}

\begin{enumerate}
\def\labelenumi{\arabic{enumi}.}
\tightlist
\item
  \textbf{Un'altalena}: va avanti e indietro, avanti e indietro\ldots{}
\item
  \textbf{Una molla}~che oscilla: su e giù, su e giù\ldots{}
\item
  \textbf{Il suono di una chitarra}: l'aria vibra in modo ripetitivo
\item
  \textbf{La luce}: anche se non la vediamo oscillare
\end{enumerate}

\hypertarget{a.-ampiezza}{%
\subsubsection{\texorpdfstring{\textbf{A. AMPIEZZA}}{A. AMPIEZZA}}\label{a.-ampiezza}}

💡

\textbf{AMPIEZZA = QUANTO E' ALTA UN'ONDA}

\begin{itemize}
\tightlist
\item
  \textbf{A parole}: quanto è ``alta'' l'onda
\item
  \textbf{Esempio concreto}:

  \begin{itemize}
  \tightlist
  \item
    Un'\textbf{onda alta 2 metri}~in mare → ampiezza grande
  \item
    Un'\textbf{onda alta 10 cm}~in una pozzanghera → ampiezza piccola
  \item
    Nell'altalena: quanto in alto va l'altalena
  \end{itemize}
\end{itemize}

\hypertarget{b.-periodo}{%
\subsubsection{\texorpdfstring{\textbf{B. PERIODO}}{B. PERIODO}}\label{b.-periodo}}

💡

\textbf{PERIODO = QUANTO IMPIEGA A FARE IL SUO CICLO COMPLETO}

\begin{itemize}
\tightlist
\item
  \textbf{A parole}: quanto tempo ci mette a fare un ciclo COMPLETO
\item
  \textbf{Esempio concreto}:

  \begin{itemize}
  \tightlist
  \item
    Se un'altalena impiega~\textbf{3 secondi}~per: andare avanti → tornare indietro → tornare al punto di partenza
  \item
    Allora il suo~\textbf{periodo è 3 secondi}
  \end{itemize}
\end{itemize}

\hypertarget{c.-frequenza}{%
\subsubsection{\texorpdfstring{\textbf{C. FREQUENZA}}{C. FREQUENZA}}\label{c.-frequenza}}

💡

\textbf{FREQUENZA = QUANTE VOLTE SI RIPETE IL CICLO IN UN DATO TEMPO}

\begin{itemize}
\tightlist
\item
  \textbf{A parole}: quante volte si ripete in un secondo
\item
  \textbf{Esempio concreto}:

  \begin{itemize}
  \tightlist
  \item
    Se il tuo cuore batte~\textbf{60 volte al minuto}
  \item
    Allora batte~\textbf{1 volta al secondo}~(frequenza = 1 battito/secondo)
  \item
    Per le onde: se un'onda si ripete 2 volte al secondo → frequenza = 2
  \end{itemize}
\end{itemize}

\hypertarget{seno-e-coseno}{%
\subsection{3. SENO E COSENO}\label{seno-e-coseno}}

💡

\textbf{SENO = FUNZIONE CHE DESCRIVE IL MOVIMENTO SU E GIU DI UN'ONDA}

\textbf{Immagina una ruota di bicicletta con un raggio rosso:}

\begin{enumerate}
\def\labelenumi{\arabic{enumi}.}
\tightlist
\item
  La ruota gira in senso antiorario
\item
  Metti un punto ROSSO sul bordo della ruota
\item
  Ora guarda SOLO l'altezza (su/giù) del punto rosso
\item
  La ruota gira → il punto va su e giù
\end{enumerate}

\textbf{Questo movimento SU e GIÙ del punto rosso è proprio la funzione SENO!}

\hypertarget{visualizziamolo-in-modo-ancora-piuxf9-concreto}{%
\subsection{\texorpdfstring{\textbf{2. VISUALIZZIAMOLO IN MODO ANCORA PIÙ CONCRETO:}}{2. VISUALIZZIAMOLO IN MODO ANCORA PIÙ CONCRETO:}}\label{visualizziamolo-in-modo-ancora-piuxf9-concreto}}

Prendi un~\textbf{compasso}:

\begin{enumerate}
\def\labelenumi{\arabic{enumi}.}
\tightlist
\item
  Apri il compasso di 5 cm (raggio)
\item
  Disegna un cerchio
\item
  Ora immagina che il compasso giri
\item
  La~\textbf{puntina che scrive}~va su e giù mentre gira
\end{enumerate}

\textbf{L'altezza della puntina rispetto al centro}~=~\textbf{VALORE DEL SENO}

\hypertarget{cosa-misura-realmente-il-seno}{%
\subsection{\texorpdfstring{\textbf{3. COSA MISURA REALMENTE IL SENO?}}{3. COSA MISURA REALMENTE IL SENO?}}\label{cosa-misura-realmente-il-seno}}

Quando hai un cerchio di raggio 1:

\begin{itemize}
\tightlist
\item
  \textbf{Seno dell'angolo}~=~\textbf{quanto sei in ALTO}~rispetto al centro
\item
  \textbf{Coseno dell'angolo}~=~\textbf{quanto sei a DESTRA}~rispetto al centro
\end{itemize}

\textbf{Esempio pratico:}

\begin{itemize}
\tightlist
\item
  Se stai in cima alla ruota panoramica (angolo = 90°):

  \begin{itemize}
  \tightlist
  \item
    Sei tutto in ALTO →~\textbf{seno = 1}~(massimo)
  \item
    Non sei né a destra né a sinistra →~\textbf{coseno = 0}
  \end{itemize}
\end{itemize}

\hypertarget{perchuxe9-ci-serve-per-le-onde}{%
\subsection{\texorpdfstring{\textbf{4. PERCHÉ CI SERVE PER LE ONDE?}}{4. PERCHÉ CI SERVE PER LE ONDE?}}\label{perchuxe9-ci-serve-per-le-onde}}

Perché il movimento~\textbf{SU e GIÙ}~della ruota che gira\ldots{}

\ldots è~\textbf{IDENTICO}~al movimento~\textbf{SU e GIÙ}~di:

\begin{itemize}
\tightlist
\item
  un'onda del mare
\item
  un'altalena
\item
  una corda di chitarra che vibra
\end{itemize}

Il~\textbf{seno}~è la~\textbf{matematica delle cose che oscillano}!

\hypertarget{gradi-e-radianti}{%
\subsection{4. GRADI E RADIANTI}\label{gradi-e-radianti}}

\hypertarget{i-gradi-li-conosci-giuxe0}{%
\subsubsection{\texorpdfstring{\textbf{I GRADI li conosci già:}}{I GRADI li conosci già:}}\label{i-gradi-li-conosci-giuxe0}}

\begin{itemize}
\tightlist
\item
  Cerchio completo = 360°
\item
  Mezzo giro = 180°
\item
  Quarto di giro = 90°
\end{itemize}

\textbf{Problema con i gradi:}~sono un'unità ``inventata'' (i Babilonesi usavano base 60)

\hypertarget{i-radianti-sono-piuxf9-naturali}{%
\subsubsection{\texorpdfstring{\textbf{I RADIANTI sono più ``naturali'':}}{I RADIANTI sono più ``naturali'':}}\label{i-radianti-sono-piuxf9-naturali}}

\textbf{Immagina di srotolare una circonferenza:}

\begin{enumerate}
\def\labelenumi{\arabic{enumi}.}
\tightlist
\item
  Prendi un cerchio
\item
  Taglia il bordo
\item
  Stendilo dritto
\item
  \textbf{La lunghezza del bordo srotolato}~si chiama~\textbf{circonferenza}
\end{enumerate}

Per un cerchio di raggio 1:

\begin{itemize}
\tightlist
\item
  Circonferenza = 2π × raggio = 2π × 1 =~\textbf{2π}
\end{itemize}

\hypertarget{cosa-sono-i-radianti-concretamente}{%
\subsection{\texorpdfstring{\textbf{2. COSA SONO I RADIANTI? CONCRETAMENTE}}{2. COSA SONO I RADIANTI? CONCRETAMENTE}}\label{cosa-sono-i-radianti-concretamente}}

\textbf{Un radiante}~= l'angolo in cui~\textbf{l'arco di circonferenza è LUNGO QUANTO IL RAGGIO}

\textbf{Esempio visivo:}

\begin{itemize}
\tightlist
\item
  Prendi una fetta di torta
\item
  Se il bordo curvo è~\textbf{lungo quanto il raggio}\ldots{}
\item
  \ldots allora l'angolo è~\textbf{1 radiante}
\end{itemize}

\textbf{Quanto è grande 1 radiante in gradi?}

\begin{itemize}
\tightlist
\item
  Cerchio completo = 360° = 2π radianti
\item
  Quindi: 1 radiante ≈ 360°/(2×3.14) ≈ 360°/6.28 ≈~\textbf{57.3°}
\end{itemize}

\hypertarget{perchuxe9-i-matematici-usano-i-radianti}{%
\subsection{\texorpdfstring{\textbf{3. PERCHÉ I MATEMATICI USANO I RADIANTI?}}{3. PERCHÉ I MATEMATICI USANO I RADIANTI?}}\label{perchuxe9-i-matematici-usano-i-radianti}}

Perché con i radianti le formule diventano più SEMPLICI:

Con gradi:

text

\begin{verbatim}
sin(90°) = 1
(sin x)' = (π/180) cos x  ← formula complicata!
\end{verbatim}

Con radianti:

text

\begin{verbatim}
sin(π/2) = 1
(sin x)' = cos x  ← formula bellissima e semplice!
\end{verbatim}

\hypertarget{conversione-facile}{%
\subsection{\texorpdfstring{\textbf{4. CONVERSIONE FACILE:}}{4. CONVERSIONE FACILE:}}\label{conversione-facile}}

\textbf{Da gradi a radianti:}

text

\begin{verbatim}
radianti = gradi × π / 180
\end{verbatim}

\textbf{Esempio pratico:}

\begin{itemize}
\tightlist
\item
  90° → 90 × π/180 = π/2 radianti
\item
  180° → 180 × π/180 = π radianti
\item
  360° → 360 × π/180 = 2π radianti
\end{itemize}

Immagina una~\textbf{torta perfettamente rotonda}:

\textbf{Situazione A - Con i GRADI:}

\begin{itemize}
\tightlist
\item
  Tagli la torta in 360 fette uguali
\item
  Ogni fetta ha angolo =~\textbf{1°}
\item
  Ma\ldots{} perché proprio 360? È un numero ``storico'' (i Babilonesi contavano in base 60)
\end{itemize}

\textbf{Situazione B - Con i RADIANTI:}

\begin{itemize}
\tightlist
\item
  Non conti le fette, ma misuri il~\textbf{bordo della torta}
\item
  Usi il~\textbf{raggio della torta}~come unità di misura
\end{itemize}

\hypertarget{esperimento-mentale-concreto}{%
\subsection{\texorpdfstring{\textbf{2. ESPERIMENTO MENTALE CONCRETO}}{2. ESPERIMENTO MENTALE CONCRETO}}\label{esperimento-mentale-concreto}}

\textbf{Materiale:}

\begin{enumerate}
\def\labelenumi{\arabic{enumi}.}
\tightlist
\item
  Una moneta da 1 euro (raggio ≈ 1.2 cm)
\item
  Un filo
\item
  Forbici
\end{enumerate}

\textbf{Procedimento:}

\begin{enumerate}
\def\labelenumi{\arabic{enumi}.}
\tightlist
\item
  Misura il~\textbf{raggio}~della moneta: ≈ 1.2 cm
\item
  Prendi il filo e tagliane un pezzo~\textbf{lungo come il raggio}~(1.2 cm)
\item
  Metti il filo sul bordo della moneta
\item
  \textbf{L'angolo che copre quel filo}~=~\textbf{1 RADIANTE}
\end{enumerate}

\textbf{Cosa succede?}

\begin{itemize}
\tightlist
\item
  Il filo (lungo come il raggio) si piega lungo il bordo
\item
  Copre un certo angolo
\item
  Quell'angolo è circa~\textbf{57.3°}
\end{itemize}

\hypertarget{perchuxe9-proprio-il-raggio}{%
\subsection{\texorpdfstring{\textbf{3. PERCHÉ PROPRIO IL RAGGIO?}}{3. PERCHÉ PROPRIO IL RAGGIO?}}\label{perchuxe9-proprio-il-raggio}}

Perché è~\textbf{NATURALE}:

Se prendo una circonferenza e la~\textbf{srotolo}:

\begin{itemize}
\tightlist
\item
  Lunghezza totale =~\textbf{2π × raggio}
\item
  Se misuro gli angoli in radianti:

  \begin{itemize}
  \tightlist
  \item
    Angolo pieno = lunghezza totale / raggio =~\textbf{2π}
  \item
    Mezzo giro = π
  \item
    Quarto di giro = π/2
  \end{itemize}
\end{itemize}

\textbf{Vantaggio:}~Le formule non hanno ``fattori strani'' come π/180!

\hypertarget{visualizziamo-con-un-cerchio-di-raggio-1}{%
\subsection{\texorpdfstring{\textbf{4. VISUALIZZIAMO CON UN CERCHIO DI RAGGIO 1}}{4. VISUALIZZIAMO CON UN CERCHIO DI RAGGIO 1}}\label{visualizziamo-con-un-cerchio-di-raggio-1}}

Prendiamo un cerchio di~\textbf{raggio = 1 unità}~(un ``cerchio unitario''):

\textbf{Con i GRADI:}

\begin{itemize}
\tightlist
\item
  360 fette
\item
  Ogni fetta: angolo 1°, arco piccolo
\end{itemize}

\textbf{Con i RADIANTI:}

\begin{itemize}
\tightlist
\item
  Circonferenza totale = 2π × 1 =~\textbf{2π unità}
\item
  Se prendo un arco lungo~\textbf{1 unità}~(come il raggio!)\ldots{}
\item
  \ldots l'angolo corrispondente =~\textbf{1 radiante}
\end{itemize}

\hypertarget{tabella-di-confronto-concreta}{%
\subsection{\texorpdfstring{\textbf{5. TABELLA DI CONFRONTO CONCRETA}}{5. TABELLA DI CONFRONTO CONCRETA}}\label{tabella-di-confronto-concreta}}

\begin{longtable}[]{@{}llll@{}}
\toprule\noalign{}
\textbf{Angolo} & \textbf{Gradi} & \textbf{Radianti} & \textbf{Cosa vedi nel cerchio} \\
\midrule\noalign{}
\endhead
\bottomrule\noalign{}
\endlastfoot
Giro completo & 360° & 2π ≈ 6.28 & Tutt'intorno \\
Mezzo giro & 180° & π ≈ 3.14 & Da sinistra a destra \\
Quarto di giro & 90° & π/2 ≈ 1.57 & Da centro a cima \\
Un ottavo & 45° & π/4 ≈ 0.785 & Metà della salita \\
\end{longtable}

\hypertarget{esercizio-pratico-guidato}{%
\subsection{\texorpdfstring{\textbf{6. ESERCIZIO PRATICO GUIDATO}}{6. ESERCIZIO PRATICO GUIDATO}}\label{esercizio-pratico-guidato}}

\textbf{Problema:}~Ho un cerchio di raggio 5 cm.

Quanto è lungo l'arco che corrisponde a π/3 radianti?

\textbf{Ragioniamo insieme:}

\begin{enumerate}
\def\labelenumi{\arabic{enumi}.}
\tightlist
\item
  Ricorda:~\textbf{1 radiante}~= arco lungo come il raggio
\item
  Qui raggio = 5 cm → 1 radiante = arco di 5 cm
\item
  L'angolo è π/3 radianti ≈ 1.047 radianti
\item
  Arco = angolo × raggio = (π/3) × 5 cm ≈ 5.24 cm
\end{enumerate}

\textbf{Formula magica:}~\texttt{Arco\ =\ Angolo(in\ radianti)\ ×\ Raggio}

Con i gradi sarebbe:~\texttt{Arco\ =\ (gradi\ ×\ π/180)\ ×\ Raggio}~→ più complicato!

\begin{enumerate}
\def\labelenumi{\arabic{enumi}.}
\tightlist
\item
  \textbf{Seno/Coseno}~→ descrivono cose che oscillano (altalene, onde)
\item
  \textbf{Radianti}~→ modo ``naturale'' per misurare angoli (usa il raggio)
\item
  \textbf{Onda sinusoidale}~= A × sin(2πft + φ) → descrive un'oscillazione pura
\item
  \textbf{Fourier}~→ dice che le cose complesse sono somme di cose semplici
\end{enumerate}

\begin{verbatim}
Cose che oscillano → Seno/Coseno → Onde pure → Somma di onde → Fourier
       ↑                    ↑              ↑            ↑           ↑
   altalene           radianti        frequenza    musica      applicazioni
   suono              formule         ampiezza     immagini    reali
   luce               semplici        fase         telecom
\end{verbatim}

\hypertarget{fourier}{%
\subsection{FOURIER}\label{fourier}}

\hypertarget{intuizione-prima-delle-formule}{%
\subsection{\texorpdfstring{\textbf{1. INTUIZIONE PRIMA DELLE FORMULE}}{1. INTUIZIONE PRIMA DELLE FORMULE}}\label{intuizione-prima-delle-formule}}

\textbf{Immagina di essere un DJ o un tecnico del suono:}

\begin{itemize}
\tightlist
\item
  Hai una~\textbf{canzone complessa}~(voce, chitarra, batteria, basso)
\item
  Ogni strumento fa~\textbf{onde diverse}
\item
  Il suono finale è un~\textbf{MISCUGLIO}~di tutte queste onde
\end{itemize}

\textbf{Domanda:}~Come fai a separare la voce dalla chitarra?

\textbf{Risposta:}~Con l'analisi di Fourier!

\hypertarget{lidea-geniale-di-fourier-1807}{%
\subsection{\texorpdfstring{\textbf{2. L'IDEA GENIALE DI FOURIER (1807)}}{2. L'IDEA GENIALE DI FOURIER (1807)}}\label{lidea-geniale-di-fourier-1807}}

Joseph Fourier scoprì che:

\begin{quote}
QUALSIASI~suono/seguale periodico, per quanto complesso,

può essere scomposto in una~\textbf{SOMMA di onde sinusoidali SEMPLICI}
\end{quote}

\textbf{Esempio concreto - Un accordo di pianoforte:}

\begin{itemize}
\tightlist
\item
  Premi DO, MI, SOL insieme
\item
  Ottieni un suono ``ricco''
\item
  Fourier dice: ``Questo suono = Onda DO pura + Onda MI pura + Onda SOL pura''
\end{itemize}

\hypertarget{visualizziamolo-con-le-onde-del-mare}{%
\subsection{\texorpdfstring{\textbf{3. VISUALIZZIAMOLO CON LE ONDE DEL MARE}}{3. VISUALIZZIAMOLO CON LE ONDE DEL MARE}}\label{visualizziamolo-con-le-onde-del-mare}}

\textbf{Scenario:}

\begin{enumerate}
\def\labelenumi{\arabic{enumi}.}
\tightlist
\item
  \textbf{Onda A}: alta 1m, periodo 5 secondi
\item
  \textbf{Onda B}: alta 0.5m, periodo 3 secondi
\item
  \textbf{Onda C}: alta 0.3m, periodo 8 secondi
\end{enumerate}

\textbf{Cosa vedi in mare?}

\begin{itemize}
\tightlist
\item
  Non vedi le 3 onde separate!
\item
  Vedi un'onda~\textbf{COMPLESSA}~che sembra irregolare
\item
  Ma in realtà è la~\textbf{SOMMA}~delle 3 onde semplici!
\end{itemize}

\hypertarget{la-serie-di-fourier-in-poche-parole}{%
\subsection{\texorpdfstring{\textbf{4. LA SERIE DI FOURIER IN POCHE PAROLE}}{4. LA SERIE DI FOURIER IN POCHE PAROLE}}\label{la-serie-di-fourier-in-poche-parole}}

\begin{quote}
Data un'onda complessa periodica,

la Serie di Fourier ci dice:

\begin{enumerate}
\def\labelenumi{\arabic{enumi}.}
\tightlist
\item
  \textbf{Quali}~onde sinusoidali ci sono dentro
\item
  Con che~\textbf{ampiezza}~(quanto sono forti)
\item
  Con che~\textbf{fase}~(quando iniziano)
\end{enumerate}
\end{quote}

\textbf{Formula concettuale (senza panico!):}

text

\begin{verbatim}
Onda complessa = A₁·sin(ω₁t+φ₁) + A₂·sin(ω₂t+φ₂) + A₃·sin(ω₃t+φ₃) + ...
\end{verbatim}

Dove ω₁, ω₂, ω₃ sono frequenze multiple della fondamentale.

\hypertarget{applicazioni-reali---dove-serve}{%
\subsection{\texorpdfstring{\textbf{5. APPLICAZIONI REALI - DOVE SERVE?}}{5. APPLICAZIONI REALI - DOVE SERVE?}}\label{applicazioni-reali---dove-serve}}

\begin{enumerate}
\def\labelenumi{\arabic{enumi}.}
\tightlist
\item
  \textbf{🎵 AUDIO DIGITALE}~(MP3, Spotify)

  \begin{itemize}
  \tightlist
  \item
    Comprimi la musica trovando le onde ``principali''
  \item
    Elimina quelle che l'orecchio non sente
  \end{itemize}
\item
  \textbf{📡 COMUNICAZIONI}~(Wi-Fi, telefoni)

  \begin{itemize}
  \tightlist
  \item
    Ogni canale ha una sua frequenza
  \item
    Fourier separa i canali
  \end{itemize}
\item
  \textbf{🩺 MEDICINA}~(ECG, risonanza magnetica)

  \begin{itemize}
  \tightlist
  \item
    L'ECG è un'onda complessa
  \item
    Fourier separa i battiti dai disturbi
  \end{itemize}
\item
  \textbf{🖼️ IMMAGINI DIGITALI}~(JPEG, Instagram)

  \begin{itemize}
  \tightlist
  \item
    Un'immagine = onde di colore in 2D
  \item
    JPEG comprime con Fourier 2D
  \end{itemize}
\end{enumerate}

\end{document}
